\section{Source corrections, verification, and further research}
\label{audit:sec}

\subsection{Corrections proposed against the pinned manuscript}

The inspection was targeted: the relevant harmonic, Dougall, twist,
collision, and Gaussian-status material was checked, rather than every
chapter and every archived computational receipt. The following two
edits are proposed for \src{chapters/03-algebraic.tex} in the
canonical manuscript \cite{RepoMain}. They concern the exact snapshot
specified in Section~\ref{scope:sec}. The package contains a small
unified patch; the repository itself has not been changed.

\paragraph{A real logarithm on the negative branch.}
In the equation labeled \src{cleo:eq:lintanh}, put
$\theta=\arcsin\lambda$ for real $|\lambda|<1$. The formula valid
on both sides of zero is
\begin{equation}\label{audit:real-atanh}
 \int_0^{\pi/2}\operatorname{arctanh}(\lambda\sin x)\,dx
 =\theta\log\left|\tan\frac\theta2\right|
       +2\operatorname{Cl}_2(\theta)
       -\frac12\operatorname{Cl}_2(2\theta).
\end{equation}
Here $\operatorname{Cl}_2(\theta)=\sum_{n\ge1}\sin(n\theta)/n^2$.
The value at $\theta=0$ is understood by continuity. The original
display omits the absolute value. For negative $\lambda$, the
real logarithm of $\tan(\theta/2)$ is not defined, whereas the
integral is real. Taking a principal complex logarithm would add
an unwanted imaginary term. This is a carry-forward correction
already identified in the incoming audit material, retained here
because the pinned canonical display still needs it.

For a direct verification, differentiate the left side after
$\lambda=\sin\theta$ to obtain $\theta/\sin\theta$.
The derivative identity
$\operatorname{Cl}_2'(\theta)=-\log|2\sin(\theta/2)|$
shows that the three logarithmic terms in the derivative of the
right side cancel, leaving the same $\theta/\sin\theta$.
Both sides tend to zero at the origin. This proves
\eqref{audit:real-atanh} on the two intervals and through their
continuous joining point.

\paragraph{Remove an unsupported non-elementarity assertion.}
The same chapter describes the $(1,1)$ diagonal as the
``simplest non-elementary diagonal.'' Its exact evaluation
\[
 \int_0^{\pi/2}\arctan^2(\sin x)\,dx
       =\pi\chi_2(3-2\sqrt2)
\]
where $\chi_2(z)=\sum_{n\ge0}z^{2n+1}/(2n+1)^2$,
does not prove non-elementarity or minimality in any specified
constant field. The chapter's later paragraph correctly says that
these properties remain unproved. We propose the neutral wording
``For the $(1,1)$ diagonal, the formula gives \dots''.
The value is unchanged. A finite unsuccessful relation search
cannot justify a stronger assertion.

\subsection{A corrected collision expectation, not a failed source theorem}

The separated-grid theorem in \cite{RepoExact} assumes disjoint
singular supports and is not contradicted by this continuation.
Nor is the proved bilinear coincidence comparison invalidated.
What fails is an unqualified extension of the positive matching
expectation in its Question~1 to geometric triple collisions.
The exact counterexample is \eqref{cl:finite-anomaly}.
A corrected research statement is:

\begin{quote}
For a specified collision path and ambient cutoff, determine the
complete local collision expression, including its finite
rate-dependent term. Subtract that expression before comparing
the limit with the merged ambient finite part.
\end{quote}

For the zero-index triple with arbitrary fixed rates $0<a<b$,
Theorem~\ref{cl:collision} proves precisely this statement, with
the complete expression
$\epsilon^{-2}P_{-2}+\epsilon^{-1}P_{-1}+P_0$.
The all-factor theorem fixes the local meromorphic completion at
a prescribed collision configuration; it does not by itself
compute every geometric collision expansion for arbitrary
Stieltjes indices and rates depending on different powers of
$\epsilon$. Those extensions are questions below.

\subsection{Statements that must retain their current status}

The source's $S_2$ and $S_4$ reductions already have proofs.
The current $S_6$ candidate is labeled
\src{cycloquot:conj:S6} in
\src{chapters/04-cyclotomic-quotients.tex}; the retained $S_8$
candidate is labeled \src{s8new:conj:S8} in
\src{chapters/04-S8-candidate.tex} \cite{RepoMain}.
Both remain conjectural here. The later $S_8$ vector must not
be confused with the different older vector that the manuscript
rigorously rejected. Exact equivalence of two candidate coordinate
vectors proves equivalence, not vanishing of either residual.
Likewise, an interval containing zero proves proximity only.
No new Gaussian integer relation was fitted in this work.

The harmonic Newton theorem solves the explicit-representation
problem for the entire remainder. It leaves the finite-basis
reduction question open at general depth and spectral order.
The Dougall Bell generator evaluates all its mixed combinations;
it does not show that those combinations span every conceivable
polygamma product. The unequal-twist scalar theorem concerns
identities of functions with two distinct branch centers;
it does not prove independence of their values on the integer
Fourier lattice. These qualifications are part of the results.

\subsection{Reproducibility and the role of computation}

The analytic proofs supply the infinite summations, domains,
continuations, and interchanges of limits. Four scripts implement
independent checks or exact finite algebra. Their JSON records
are included in \src{results/}; they contain the actual inputs,
working precision, and residuals. The scripts need Python,
\src{mpmath}, and \src{sympy}. The exact environment versions
and commands are recorded in the package's \src{README.txt}.

\begin{longtable}{@{}>{\raggedright\arraybackslash}p{0.28\textwidth}>{\raggedright\arraybackslash}p{0.63\textwidth}@{}}
\toprule
Artifact & Checks and interpretation\\
\midrule
\endhead
\src{verify_harmonic.py}
& Twenty-six numerical checks: twelve truncated Newton spectral
derivatives against Hurwitz derivatives, one complex-parameter
Mellin comparison, nine finite integer-shift specializations,
and four repeated antiderivatives. The Newton series uses
$240$ terms and $180$ working digits; its observed discrepancies
are about $10^{-18}$ to $10^{-15}$, set by truncation rather
than the working precision.\\[6pt]
\src{verify_dougall.py}
& Thirteen numerical comparisons of the trigamma square,
two-parameter Gamma identity, mixed spectral generator, and first
spectral derivative. Six exact Bell-polynomial checks and
twenty-one exact residue-grid factorizations use rational
symbolic arithmetic. The numerical residuals are below
$5.4\times10^{-46}$ in the supplied run.\\[6pt]
\src{verify_twists.py}
& Forty-nine exact partial-fraction checks through $p,q=7$,
one deliberately corrupted-sign rejection, and the exact
lowest-jet polarization identity. Eighteen numerical comparisons
cover ordinary Green convolutions, confluent twists, the bounded
logarithm difference, rational digamma conversion, branch
crossings, Fourier coefficients, and complex-order tail modes.\\[6pt]
\src{verify_collisions.py}
& Seventeen separated-collision diagnostics at $70$ working
digits, plus a two-split evaluation of the merged integral.
The independent interval quadratures visibly subtract each
simple pole. Their residuals test the predicted asymptotic
expansion; they are not exact-zero tests.\\
\bottomrule
\end{longtable}

The finite symbolic checks do not formalize the analytic proofs.
The floating-point comparisons are not interval enclosures.
In particular, an estimated asymptotic tail term is not a
certified error bound for the Gamma-ratio computations.
The harmonic finite differences can lose many digits through
cancellation; a larger working precision alone does not remove
the separate Newton truncation error. These are concrete
limitations of the supplied calculations, not qualifications
of the proved normal-convergence theorems.

\subsection{Further questions with concrete first targets}

The following are proposals for additional exact research. They
are not included in the proved theorem list.

\begin{question}[Which individual polygamma moments can be separated?]
At a fixed total order $w=r+s$, determine the exact span of the
mixed Bell polynomials $\mathcal B_{rs}$, allowing the elementary
linear polygamma rows and explicit parameter symmetries. The
square identity used the difference of two rows at $w=4$.
At $w=6$ and $w=8$, give exact rational elimination matrices and
identify which individual moments remain outside that span.
An absence from this specified span should be reported as such,
not as an impossibility of evaluation by other identities.
\end{question}

\begin{question}[Compact formulas at every integer residue intersection]
Use \eqref{dg:grid-first-mixed} to produce short explicit
tables at $(b,c)=(1,2)$ and $(2,2)$ for $N=2,3$.
Can central factorial polynomials organize the subtraction
coefficients without expanding large Bernoulli polynomials?
A useful certificate would retain the finite small-index
zeros and replay the mixed derivative using only exact
polynomial operations before evaluating polygamma endpoints.
\end{question}

\begin{question}[Finite reductions of regular harmonic spectral jets]
The Newton formula supplies an exact convergent representation;
classify parameter configurations at which it reduces to finitely
many ordinary Hurwitz-zeta order derivatives. Positive integral
$a$ and nonnegative integral $u$ give starting cases. For
positive depth and nonintegral rational $a$, first examine the
first two derivatives at $s=0,-1$. State the permitted function
class before asking whether a multiple Arakawa--Kaneko order
derivative is needed.
\end{question}

\begin{question}[The exact boundary of shift continuation]
The proved depth domain is $\Re a>-r$, and the Taylor series
about $a=1$ converges at least for $|a-1|<r+1$.
Determine the actual singular germ at $a=-r$ for generic complex
$s$ and compare it with the exceptional nonpositive integral
$s$, where the finite Stirling formula is polynomial in $a$.
This should distinguish a sufficient convergence domain from
a maximal analytic continuation and explain any exceptional
cancellation of the first shift singularity.
\end{question}

\begin{question}[Rational-shift character cancellations]
For rational $a$, transform the full mixed spectral germ
\eqref{dg:all-Stieltjes} into finite residue-class or
Dirichlet-character coordinates. The first target is $a=1/2$,
where $\gamma_m(a/2)=\gamma_m(1/4)$ appears. At each prescribed
order, compute an exact reduction matrix including every
conductor logarithm. Which parameter symmetries cancel a
whole character sector before numerical evaluation?
\end{question}

\begin{question}[Unequal twists approaching an integer resonance]
The finite-mode decomposition in \eqref{tw:tailseries} gives
the exact modes that change phase. Let one or both twists
approach an integer while the spectral orders are also
differentiated. After isolating the resonant Fourier modes,
determine a holomorphic completed germ and its endpoint contact
terms. The first target is a simultaneous pair approaching
zero at distinct fixed rates. The normalization should be
compared with the already known single-twist degeneration.
\end{question}

\begin{question}[From scalar functional reduction to sampled identities]
The finite scalar criterion proved here uses local monodromy.
Which additional hypotheses would let one infer a scalar
identity from equality on every integer Fourier mode in the
same finite power--logarithm class? An exact asymptotic uniqueness
theorem at infinity is one possible route. Conversely, seek
explicit sampled identities that do not extend as scalar
functional identities. This is distinct from any question of
arithmetic independence of individual special values.
\end{question}

\begin{question}[A complete geometric collision calculus]
Extend Theorem~\ref{cl:collision} to arbitrary Stieltjes indices
and argument-derivative orders, first for three unit-frequency
factors. The singular pieces are explicit derivatives of
$x^{-1}\log^m x$, so finite partial fractions and order
derivatives suggest an algorithm. It should output every
power--logarithm term and the finite rate-dependent term, with
an integrable uniform remainder. The desired result is an
identity for a specified normalization, not a sharper bound.
\end{question}

\begin{question}[Hierarchical collisions and composition]
Compare the paths $(0,\epsilon,\epsilon+\epsilon^2)$ and
$(0,\epsilon,2\epsilon)$ with iterated merging of a selected
pair followed by the third point. Can the finite corrections
be organized by nested subsets so that renormalized collision
composition is associative? The rate dependence already proved
shows why plain constant extraction is insufficient. Specify
the intermediate coordinate and subtraction at every stage.
\end{question}

\begin{question}[Permutation cancellations beyond total derivatives]
The completed generators distinguish global Tornheim data from
local endpoint data. Determine which nontrivial permutation
or character projections remove the global part for four
factors, after the elementary total-derivative identities have
been accounted for. A concrete target is total argument order
three with equal Stieltjes indices, retaining the complete
resonant-subset numerators. Produce either an exact cancellation
or a precise surviving multiple-sum kernel.
\end{question}

\begin{question}[Finite global kernels for four colliding grids]
Theorem~\ref{cl:allfactor} solves the local completion for any
number of factors. Construct a useful global spectral realization
for four frequencies on the rank-three lattice
$\sum_i p_i n_i=0$, with a finite sign-region and character
decomposition. The first target is $(p_1,p_2,p_3,p_4)=(1,1,2,3)$.
Local existence alone should not be presented as a minimal
finite reduction of the global multiple Tornheim kernels.
\end{question}

\begin{question}[The frozen Gaussian candidates]
Retain the exact source vectors for $S_6$ and the current $S_8$.
Seek an additional functional or iterated-integral relation
beyond the projection mechanisms that already prove $S_4$.
A successful proof should provide a finite exact row certificate
and an analytic justification for every generating row,
including its endpoint evaluation. An exclusion from a chosen
finite relation space would identify a missing mechanism,
not disprove the numerical identity.
\end{question}

\subsection{Recommended integration order}

The four mathematical sections can be integrated independently after
their cited antecedents. Preserve the stated domains, branches, and
normalizations, and review the two source edits separately.
The accompanying integration notes map labels and source targets
while preserving the existing conjecture status.
